\documentclass[11pt,a4paper]{ctexart}

\usepackage{amsmath,amssymb,mathtools}
\usepackage{geometry}
\usepackage{hyperref}
\usepackage{enumitem}
\usepackage{xcolor}

\geometry{margin=2.4cm}
\hypersetup{
  colorlinks=true,
  linkcolor=blue!60!black,
  urlcolor=blue!60!black,
  citecolor=blue!60!black
}
\setlist[itemize]{leftmargin=2em}
\setlist[enumerate]{leftmargin=2em}

\title{World Action Models: 数学化学习笔记}
\author{}
\date{\today}

\begin{document}
\maketitle

\section{学习目标}

World Action Model，简称 WAM，不应只理解为“会生成未来视频的机器人模型”。更严格地说，WAM 研究的是：
\[
  \text{在部分可观测、语言条件、连续控制的序列决策问题中，如何联合建模世界动态、动作分布与任务价值。}
\]
它处在 POMDP、world model、imitation learning、video prediction、diffusion policy 和 model-based planning 的交叉处。

\section{从 POMDP 定义问题}

具身控制任务可以形式化为一个部分可观测马尔可夫决策过程：
\[
  \mathcal{M}=(\mathcal{S},\mathcal{A},\mathcal{O},T,\Omega,R,\gamma).
\]
其中：
\begin{itemize}
  \item $\mathcal{S}$ 是真实环境状态空间，通常不可完全观测；
  \item $\mathcal{A}$ 是动作空间，例如关节角速度、末端位姿增量或夹爪开合；
  \item $\mathcal{O}$ 是观测空间，例如 RGB 图像、多视角图像和 proprioception；
  \item $T(s_{t+1}\mid s_t,a_t)$ 是环境动力学；
  \item $\Omega(o_t\mid s_t)$ 是观测模型；
  \item $R(s_t,a_t)$ 是奖励函数；
  \item $\gamma$ 是折扣因子。
\end{itemize}

实际机器人通常不能直接访问 $s_t$，只能得到观测 $o_t$。若任务由语言给定，记语言指令为 $\ell\in\mathcal{L}$。一个普通策略模型学习：
\[
  \pi_\theta(a_t\mid o_{\le t},a_{<t},\ell).
\]
很多 VLA 模型进一步近似为反应式策略：
\[
  \pi_\theta(a_t\mid o_t,\ell).
\]

\section{WAM 的核心建模对象}

WAM 的关键不是只预测动作，而是联合建模动作与未来世界。一个典型 WAM 近似如下条件分布：
\[
  p_\theta(o_{t+1:t+H},a_{t:t+H-1}\mid o_{\le t},\ell).
\]
如果引入机器人自身状态 $q_t$，例如关节角、末端位姿和夹爪状态，则可写为：
\[
  p_\theta(o_{t+1:t+H},q_{t+1:t+H},a_{t:t+H-1}
  \mid o_{\le t},q_{\le t},\ell).
\]
如果模型还预测价值，则可扩展为：
\[
  p_\theta(o_{t+1:t+H},q_{t+1:t+H},a_{t:t+H-1},V_{t+H}
  \mid o_{\le t},q_{\le t},\ell).
\]
这类建模方式对应 Cosmos Policy 一类方法：把 action、future state 和 value 都纳入统一生成框架。

\section{VLA 与 WAM 的区别}

VLA 主要学习：
\[
  p(a_t\mid o_t,\ell).
\]
WAM 则学习：
\[
  p(a_{t:t+H-1},o_{t+1:t+H}\mid o_{\le t},\ell).
\]
因此，VLA 更像是
\[
  \text{observation}\rightarrow \text{action},
\]
而 WAM 更像是
\[
  \text{observation}\rightarrow \text{future world}+\text{action}.
\]

从 representation learning 的角度看，WAM 的潜在收益来自未来预测目标对隐表示 $z_t$ 的约束。理想情况下，$z_t=f_\theta(o_{\le t})$ 需要携带物理动态、物体交互和任务进展信息。可以直观理解为提高：
\[
  I(z_t;\,o_{t+1:t+H}\mid o_{\le t},a_{t:t+H-1},\ell),
\]
即隐表示与未来观测之间的条件互信息。

\section{三类 WAM 范式}

\subsection{Joint Video-Action Modeling}

联合视频与动作建模直接学习：
\[
  p_\theta(o_{t+1:t+H},a_{t:t+H-1}\mid o_{\le t},\ell).
\]
若使用扩散模型，可令
\[
  x_0=[o_{t+1:t+H},a_{t:t+H-1}],\qquad c=[o_{\le t},\ell].
\]
标准 denoising loss 为：
\[
  \mathcal{L}_{\mathrm{diff}}
  =
  \mathbb{E}_{x_0,\epsilon,\tau}
  \left[
    \left\|
      \epsilon-\epsilon_\theta(x_\tau,\tau,c)
    \right\|_2^2
  \right].
\]
这里未来视频 token 和 action token 被放在同一个生成对象中。

\subsection{Imagine-Then-Execute}

这类方法先生成未来：
\[
  \hat{o}_{t+1:t+H}\sim
  p_\theta(o_{t+1:t+H}\mid o_{\le t},\ell),
\]
再根据想象出的未来预测动作：
\[
  a_{t:t+H-1}\sim
  \pi_\phi(a_{t:t+H-1}\mid o_{\le t},\hat{o}_{t+1:t+H},\ell).
\]
直觉是：如果模型能想象任务成功的未来状态，动作模型就可以反推出通往该未来的动作。缺点是测试时通常需要多步视频扩散去噪，推理延迟高。

\subsection{Fast-WAM}

Fast-WAM 将训练目标和推理过程解耦。训练时仍使用视频预测：
\[
  \mathcal{L}
  =
  \mathcal{L}_{\mathrm{action}}
  +
  \lambda\mathcal{L}_{\mathrm{video}}.
\]
但测试时不显式生成未来视频，而是直接预测动作：
\[
  a_{t:t+H-1}\sim \pi_\theta(a_{t:t+H-1}\mid z_t,\ell),
  \qquad z_t=f_\theta(o_{\le t}).
\]
因此，Fast-WAM 的核心问题是：WAM 的收益究竟来自测试时 future imagination，还是来自训练时 video prediction 对世界表征的塑造？

\section{阅读论文时应追问的问题}

\begin{enumerate}
  \item 模型实际近似的是哪个条件分布？
  \item 未来观测、动作、proprioception 和 value 是联合生成，还是分阶段生成？
  \item 视频预测目标只用于训练，还是也参与测试时决策？
  \item 动作是单步预测，还是 action chunk？
  \item 推理时是否进行 best-of-$N$ sampling、value ranking 或 model-based planning？
  \item 性能提升来自更强的世界表征，还是来自显式规划？
\end{enumerate}

\section{建议阅读顺序}

建议先读 Fast-WAM，因为它的问题意识最清楚：WAM 是否真的需要测试时想象未来。然后读 Cosmos Policy，因为它是更完整的系统，包含 action、future state、value 和 planning。最后再读 WM as simulator，因为它更像是把 world model 当作数据生成器或仿真器的应用，而不是 WAM 的入门核心。

\end{document}
